\documentclass{beamer}
%Information to be included in the title page:
\title{Computing psuedorandom functions within Ren}

\author{}
\institute{Field Extension Institute}

\date{RW10, 8th August 2026}

\begin{document}

\frame{\titlepage}

\begin{frame}
\frametitle{Motivation}
\begin{itemize}
    \item Sources of cryptographic (psuedo)randomness are generally useful things to have.
    \item Plaintext-generated and public-key-encrypted randomness may not be viable (eg. untrusted server or other party).
\end{itemize}

\vspace{1cm}
\textit{Specific use case}: rngeo, an FHE-based location game where scoring was based on a random number generator seeded by location and timestamp\footnote{* strictly speaking, we're using a pseudorandom number generator instead of a true PRF}, among other things.


\end{frame}

\begin{frame}
\frametitle{The problem}
Can we efficiently (in depth, compute, memory use) implement a pseudorandom function as a circuit in Ren which takes encrypted inputs?
\end{frame}

\begin{frame}
\frametitle{Common PRFs and barriers we observed}

\begin{itemize}
    \item Hash-based: hash a counter and secret seed to get a pseudorandom output. 
    Something like SHA-256, but that requires mod-$2^{32}$ addition which is tricky whether we encode inputs as ints or bitstrings (adders necessarily have to be deep because of carries).
    \item Cipher-based: either block or stream ciphers, generating a sequence of pseudorandom outputs.
    AES has the SubBytes step which is a byte-substitution lookup table, and ChaCha also requires mod-$2^{32}$ addition so neither are particularly convenient to implement.
\end{itemize}
\end{frame}

\begin{frame}
\frametitle{Some FHE-friendly PRF constructions}

\begin{itemize}
    \item Simple substitution-permutation networks (SPNs) which use FHE-friendly substitution (eg. polynomials) and permutation (eg. rotations) layers.
    This did not achieve very good confusion/diffusion metrics - Shannon's properties of secure ciphers.
    \item Rasta (\href{https://eprint.iacr.org/2018/181.pdf}{eprint 2018/181}) is a stream cipher using alternating substitution and affine transforms over a binary bitstream.
    It is designed to have low AND-depth and use few ANDs per bit, with the only other operation being XOR (both of which are relatively simple in Ren).
    \item HERA (\href{https://eprint.iacr.org/2020/1335.pdf}{eprint 2020/135}) was proposed as part of a CKKS/FV transciphering framework, and is a stream cipher on $\mathbb{Z}_t$.
    It is more performant than Rasta, but requires evaluating operations over a general modulus $t$.
\end{itemize}
\end{frame}

\begin{frame}
\frametitle{Current status}
rngeo uses a Rasta-like PRF implementation on a 256-bit input; the same substitution layer as Rasta, but a different affine layer (Rasta's uses a matrix multiplication).

\vspace{0.5cm}
\textit{Current numbers:} 256-bit input/output, 4 rounds of alternating affine/subst. followed by a final affine, 8 rotation keys (at rngeo secure params, $\approx 715$ MB per key), ~6 depth per round.
Diffusion per bit-flip $\approx 0.39$, per-slot and pooled digit distributions within 15\% of uniform. 

\vspace{0.5cm}
\textit{Potential holes:} the round constants for each affine layer are still generated server-side in plaintext. 
\end{frame}

\end{document}